%! Inverses of Exponential Functions — Unit 5, Lesson 5.2 — Homework
\documentclass[10pt]{article}
\usepackage{precalculus-article}
\usepackage{precalculus-boxes}
\newcommand{\TallMath}[1]{$\displaystyle #1\rule[-1.4em]{0pt}{3.2em}$}

\begin{document}

\pageheader{Unit 5, Lesson 5.2}{Homework}
\namedateperiod

\vspace{0.1in}

\begin{enumerate}[leftmargin=1.8em, itemsep=14pt]

\item Let $g(x) = 6^x$.

\begin{enumerate}[leftmargin=1.5em, itemsep=6pt, label=\alph*.]
  \item Write the inverse function $f(x)$.

  \vspace{4pt}
  $f(x) = $ \blank{4cm}

  \item Complete the table, then write the corresponding ordered pairs for $f$.

  \vspace{4pt}
  \begin{center}
  \renewcommand{\arraystretch}{1.8}
  \begin{tabular}{|c|c|c|}
  \hline
  \rowcolor{sky}
  $x$ & $g(x) = 6^x$ & Point on $f$: $(g(x),\; x)$ \\
  \hline
  $0$ & \blank{1.5cm} & \blank{3cm} \\
  \hline
  $1$ & \blank{1.5cm} & \blank{3cm} \\
  \hline
  $2$ & \blank{1.5cm} & \blank{3cm} \\
  \hline
  \end{tabular}
  \end{center}
\end{enumerate}

\item Verify that $f(x) = \log_5 x$ and $g(x) = 5^x$ are inverse functions by computing
      both compositions. Show all steps.

\vspace{4pt}
$f(g(x)) = \log_5(5^x) = $ \blank{3cm}

\vspace{6pt}
$g(f(x)) = 5^{\log_5 x} = $ \blank{3cm}

\vspace{4pt}
Conclusion: \writeline

\item A biologist models algae coverage with $A(t) = 10^t$ cm$^2$, where $t$ is time in days.

\begin{enumerate}[leftmargin=1.5em, itemsep=6pt, label=\alph*.]
  \item Write the inverse function $t(A)$ that gives the number of days from the area.

  \vspace{4pt}
  $t(A) = $ \blank{4cm}

  \item Evaluate $t(1000)$. What does this value represent in context?

  \vspace{4pt}
  $t(1000) = $ \blank{3cm}

  \vspace{4pt}
  In context: \writeline
\end{enumerate}

\item \textbf{Conceptual question.} Explain in your own words why the logarithmic function's
      growth pattern (output changes additively as input changes multiplicatively) is the
      ``opposite'' of the exponential function's growth pattern.

\writelines{3}

\item Let $h(x) = 2\log_3 x$.

\begin{enumerate}[leftmargin=1.5em, itemsep=6pt, label=\alph*.]
  \item Identify the base $b$ and the value of $a$.

  \vspace{4pt}
  $b = $ \blank{2cm} \qquad $a = $ \blank{2cm}

  \item Evaluate $h(9)$ and $h(27)$.

  \vspace{4pt}
  $h(9) = $ \blank{3cm} \qquad $h(27) = $ \blank{3cm}

  \item How do the outputs of $h(x)$ compare to the outputs of $\log_3 x$ for the same
        inputs? Describe the effect of the vertical dilation by $a = 2$.

  \writeline
\end{enumerate}

\item \textbf{AP-style multiple choice.} The graph of $f(x) = \log_4 x$ is a reflection of
      which of the following over the line $y = x$?

\vspace{6pt}
\begin{enumerate}[leftmargin=2em, itemsep=4pt, label=(\Alph*)]
  \item $g(x) = x^4$
  \item $g(x) = 4x$
  \item $g(x) = 4^x$
  \item $g(x) = \log x$
\end{enumerate}

\end{enumerate}

\vspace{0.12in}

\begin{reflectionbox}
What is one thing from today's lesson that you feel confident about, and one question you
still have about inverse functions or logarithms?

\writelines{2}
\end{reflectionbox}

\end{document}
